\uuid{28Gv}
\titre{Nombres premiers}
\auteur{bodin ; croizat ; sacre}
\organisation{exo7}
\licence{CC BY-NC-SA 4.0}
\datecreate{2022-03-08}
\niveau{L1}
\module{Algèbre}
\chapitre{Arithmétique dans Z}
\sousChapitre{Nombres premiers, nombres premiers entre eux}
\difficulte{1}
\theme{Nombres premiers}
\exo7qcmid{hilisit.arithmetique.028}

\contenu{
\begin{question}
\qclassification{103.04}
    Quelles sont les affirmations vraies ?
    \begin{answers}
        \bad{Tout nombre impair supérieur à $3$ est premier.}
        \good{Tout nombre premier supérieur à $3$ est impair.}
        \good{Il existe une infinité de nombres premiers impairs.}
        \bad{Il existe une infinité de nombres premiers pairs.}
    \end{answers}
    \begin{explanations}
    Par exemple $9$ est un nombre impair qui n'est pas premier.
    Le seul nombre premier pair est $2$, tous les autres sont impairs et il y en a une infinité.
    \end{explanations}
\end{question}
}
